\documentclass[12pt]{article}
\usepackage[margin=1in]{geometry}
\usepackage{setspace}
\usepackage{parskip}

\title{Peer Review: HuanXi Zhang --- Third Draft}
\author{Alejandro De La Torre}
\date{April 2026}

\begin{document}

\maketitle

\singlespacing

\noindent \textbf{Summary and Motivation.} \\
To my understanding, the paper asks whether the expansion of gig-economy platforms in Chinese cities changes how rural households allocate labor and land, and whether these changes translate into differences in agricultural outcomes. Very cool! I especially appreciate how you connect the gig economy to classic questions in development (structural transformation, land markets, and labor reallocation). The paper is ambitious in an insanely cool way, and the combination of newly constructed platform entry data with the NFP household panel is a clear strength.

\vspace{0.5em}
\noindent \textbf{What I found most convincing.} \\
I found the core empirical results on labor reallocation and land transfer to be the most convincing part of the paper. The finding that platform entry reduces agricultural participation and increases land transfer-out is intuitive and well-supported across specifications. I also thought the income and employment decomposition (local vs. migrant work) was a strong piece of evidence that helps connect your estimates to a plausible mechanism. Overall, I think the paper makes a meaningful contribution by bringing platform expansion into the study of rural factor reallocation.

\vspace{0.5em}
\noindent \textbf{Main suggestions for improvement.}

\textit{1. Clarifying the mechanism and exposure.} \\
I found myself wanting more clarity on how rural households are actually exposed to platform opportunities. Since the treatment is defined at the city level while outcomes are measured for rural households, I think it would strengthen the paper to more explicitly discuss whether the mechanism operates through direct participation (e.g., commuting to nearby urban centers) or through indirect spillovers. I wonder if you could exploit variation in distance to the urban core or provide more direct evidence linking households to platform-related work. Even reframing this as “short-distance reallocation” rather than strictly “local” could make the interpretation clearer while still preserving the contribution.

\textit{2. Treatment construction and validation.} \\
I think your data construction using the Wayback Machine is creative and impressive. At the same time, I found myself wanting more discussion of potential measurement issues. For example, larger cities may be more likely to appear earlier in archived webpages, and webpage presence may not perfectly correspond to actual platform operations. I think the paper would benefit from a clearer validation exercise (even a simple external comparison or alternative proxy, such as search intensity or recruitment activity) to build confidence in the treatment variable.

\textit{3. Framing around migration.} \\
I really like the idea that platform expansion could allow reallocation without large-scale migration. At the same time, I found the “no migration” framing slightly stronger than what the evidence seems to support. My reading of the results is that local nonfarm opportunities increase without reducing migration, rather than fully substituting for it. I think a slightly more cautious framing—emphasizing complementarity rather than substitution—would make the contribution more precise and defensible.

\textit{4. Pre-trends and identification.} \\
I thought the event-study figures were helpful, but I found myself wanting a bit more formal discussion of pre-trends. In particular, even if individual coefficients are not statistically significant, there may still be underlying trends. I think adding a joint test of pre-treatment coefficients or discussing sensitivity to pre-trend assumptions would strengthen the credibility of the design. Relatedly, I wonder if you might consider foregrounding the Callaway–Sant’Anna results more prominently, since they are well-suited to staggered adoption settings like yours.

\textit{5. Agricultural outcomes interpretation.} \\
I found the agricultural outcomes section interesting but harder to interpret. Some results—especially for land-transferring-in and non-transferee households—seem heterogeneous and somewhat noisy. I think this section could be strengthened either by providing a clearer unifying mechanism (for example, shifts toward cash crops or capital intensification) or by narrowing the focus to the most robust outcomes. As it stands, I occasionally found it difficult to map the estimates to a single economic story.

\vspace{0.5em}
\noindent \textbf{Exposition.} \\
I found the overall structure of the paper clear and easy to follow (it is very neat and organized; it reads like a published AEA paper), especially the progression from motivation to identification to results. In a few places, I felt that simplifying or tightening sentences (particularly in the introduction and abstract) could make the main ideas stand out more clearly. I also think emphasizing the central contribution earlier—what we learn here that we did not know before—would help guide the reader.

\vspace{0.5em}
\noindent \textbf{Overall.} \\
I think this is a strong and thoughtful paper with a clear and important question. My main suggestions are about sharpening the mechanism, clarifying the interpretation of the treatment, and tightening how the results are framed. With those refinements, I think the paper has the potential to make a very compelling contribution to the literature on structural transformation and the gig economy.

\end{document}
